% term_index_mwe.tex
% Requires: pdflatex -> makeindex -> pdflatex -> pdflatex
% (or: latexmk -pdf term_index_mwe.tex)
\documentclass{article}

\usepackage{imakeidx}   % index with hyperlinked page numbers
\usepackage{hyperref}   % hyperlinks; load before xcolor to avoid option clash
\usepackage{xcolor}     % \textcolor
\usepackage{xparse}     % \NewDocumentCommand for starred/optional variants

% ---- Set up the index -------------------------------------------------------
\makeindex[title=Index,intoc,columns=2]

% ---- Aux-writing helpers for \refterm ---------------------------------------
%
% \term[label]{word} writes \termdef{label}{word} to the .aux file.
% On the next pass this defines \termtext@label, which \refterm reads.
\makeatletter
% Write to aux:
\newcommand{\writetermaux}[2]{%
  \protected@write\@auxout{}{%
    \string\termdef{#1}{#2}%
  }%
}

% Read from aux (defines \termtext@<label>):
\newcommand{\termdef}[2]{%
  \global\@namedef{termtext@#1}{#2}%
}

% ---- \term command ----------------------------------------------------------
%
%   \term[label]{word}  -- blue, indexed, and labelled for \refterm
%   \term{word}         -- blue and indexed (no label)
%   \term*{word}        -- blue only (no index entry, no label)
%   \term*[label]{word} -- blue only, but still labelled (edge case)

\NewDocumentCommand{\term}{ s o m }{%
    \textcolor{blue}{#3}\IfBooleanF{#1}{\index{#3}}%
    \IfValueT{#2}{% label provided: write aux + set hyperref anchor
        \writetermaux{#2}{#3}%
        \phantomsection\label{term:#2}%
    }%
}

% ---- \refterm command -------------------------------------------------------
%
%   \refterm{label}  -- hyperlink to the \term[label]{...} occurrence,
%                       displaying the word used there in blue.
%
% Falls back to "??" (like \ref on first pass) if the label is unknown.

\newcommand{\refterm}[1]{%
    \hyperref[term:#1]{%
        \@ifundefined{termtext@#1}{\textcolor{red}{UNDEFINED TERM REFERENCE}}{\textcolor{blue}{\csuse{termtext@#1}}}%
    }%
}
\makeatother
% -----------------------------------------------------------------------------

\begin{document}

\tableofcontents
\bigskip

\section{Limits}

We begin with the notion of a \term[lim-def]{limit}.  Informally, the
\term{limit} of a function $f(x)$ as $x \to a$ is the value that
$f(x)$ approaches as $x$ gets arbitrarily close to~$a$.

\begin{quote}\textbf{Definition.}\quad
  Let $f$ be defined near $a$.  We say the \term{limit} of $f$ at $a$
  equals $L$, written $\lim_{x\to a} f(x) = L$, if for every
  $\varepsilon > 0$ there exists $\delta > 0$ such that
  $|f(x) - L| < \varepsilon$ whenever $0 < |x - a| < \delta$.
\end{quote}

A function that possesses a limit everywhere is called
\term*[ttt]{well-behaved} (informal; coloured but not indexed).

\section{Continuity}

A function is \term[cont-def]{continuous} at a point $a$ if
$\lim_{x \to a} f(x) = f(a)$.  \term{Continuity} on an interval means
\term{continuous} at every point of that interval.

We use \term{limit} once more here to emphasise multiple index entries.

\section{Back-references}

Recall the \refterm{lim-def} introduced in Section~1, and the notion
of \refterm{cont-def} defined above.

\refterm*{ttt}

\printindex

\end{document}